\documentclass[11pt,a4paper]{article}
\usepackage[utf8]{inputenc}
\usepackage{amsmath}
\usepackage{amsfonts}
\usepackage{amssymb}
\usepackage{graphicx}
\usepackage{listings}
\usepackage{xcolor}
\usepackage{geometry}
\usepackage{hyperref}
\usepackage{fancyhdr}
\usepackage{tcolorbox}

\geometry{margin=2.5cm}
\pagestyle{fancy}
\fancyhf{}
\fancyhead[L]{LSH Collision Test Documentation}
\fancyhead[R]{\today}
\fancyfoot[C]{\thepage}

% Code listing style
\lstdefinestyle{cpplstyle}{
    language=C++,
    backgroundcolor=\color{gray!10},
    commentstyle=\color{green!60!black},
    keywordstyle=\color{blue},
    numberstyle=\tiny\color{gray},
    stringstyle=\color{red},
    basicstyle=\footnotesize\ttfamily,
    breakatwhitespace=false,
    breaklines=true,
    captionpos=b,
    keepspaces=true,
    numbers=left,
    numbersep=5pt,
    showspaces=false,
    showstringspaces=false,
    showtabs=false,
    tabsize=2,
    frame=single,
    rulecolor=\color{gray!30}
}
\lstset{style=cpplstyle}

\title{\textbf{LSH Collision Test Documentation}\\
\large{Locality-Sensitive Hashing Collision Analysis for SMHasher3}}
\author{SMHasher3 Project}
\date{\today}

\begin{document}

\maketitle

\tableofcontents
\newpage

\section{Overview}

The \texttt{LSHCollisionTest.cpp} file implements a comprehensive testing framework for evaluating the locality-sensitive hashing (LSH) properties of hash functions within the SMHasher3 test suite. This test specifically focuses on analyzing collision rates between similar biological sequences (DNA sequences) to assess how well a hash function preserves locality in the input space.

\subsection{Purpose}

The primary purpose of this test is to evaluate whether hash functions exhibit appropriate locality-sensitive properties when applied to biological sequence data. In an ideal LSH scenario:
\begin{itemize}
    \item Similar input sequences should have a higher probability of hash collisions
    \item Dissimilar sequences should have a lower probability of hash collisions
    \item The collision probability should correlate with sequence similarity
\end{itemize}

\section{Mathematical Foundation}

\subsection{DNA Sequence Representation}

DNA sequences are represented using a 2-bit encoding scheme:
\begin{align}
\text{Base Encoding:} \quad A \rightarrow 00_2, \quad C \rightarrow 01_2, \quad G \rightarrow 10_2, \quad T \rightarrow 11_2
\end{align}

For a sequence of length $n$ bases, the total storage requirement is:
\begin{equation}
\text{Storage Bytes} = \left\lceil \frac{2n}{8} \right\rceil = \left\lceil \frac{n}{4} \right\rceil
\end{equation}

\subsection{Similarity Metrics}

The test uses sequence similarity based on Hamming distance. For two sequences $S_1$ and $S_2$ of equal length $n$:

\begin{equation}
\text{Hamming Distance}(S_1, S_2) = \sum_{i=1}^{n} [S_1[i] \neq S_2[i]]
\end{equation}

\begin{equation}
\text{Similarity}(S_1, S_2) = 1 - \frac{\text{Hamming Distance}(S_1, S_2)}{n}
\end{equation}

\subsection{Collision Rate Analysis}

For a hash family $\mathcal{H}$ and two sequences $S_1, S_2$, the collision probability is:

\begin{equation}
P_{collision}(S_1, S_2) = \frac{1}{|\mathcal{H}|} \sum_{h \in \mathcal{H}} [h(S_1) = h(S_2)]
\end{equation}

\section{Implementation Details}

\subsection{Test Configuration}

\begin{tcolorbox}[colback=blue!5!white,colframe=blue!75!black,title=Key Parameters]
\begin{itemize}
    \item \textbf{Key Count}: 1024 sequences per similarity level
    \item \textbf{Hash Count}: 1024 hash functions per sequence pair
    \item \textbf{Key Sizes}: Multiple bit lengths (32, 48, 64, 96, 128, 160*, 256*)
    \item \textbf{Similarity Ranges}: 0.0 to 1.0 in 0.1 increments
    \item \textbf{Output}: CSV file with collision rates
\end{itemize}
\small{*Available with \texttt{--extra} flag for non-slow hash functions}
\end{tcolorbox}

\subsection{Key Size to Sequence Length Mapping}

The relationship between key bits and DNA sequence length is:
\begin{equation}
\text{Sequence Length} = \frac{\text{Key Bits}}{2}
\end{equation}

\begin{center}
\begin{tabular}{|c|c|c|}
\hline
\textbf{Key Bits} & \textbf{Key Bytes} & \textbf{Sequence Length (bases)} \\
\hline
32 & 4 & 16 \\
48 & 6 & 24 \\
64 & 8 & 32 \\
96 & 12 & 48 \\
128 & 16 & 64 \\
160 & 20 & 80 \\
256 & 32 & 128 \\
\hline
\end{tabular}
\end{center}

\section{Algorithm Workflow}

\subsection{Main Test Function}

The \texttt{LSHCollisionTest} function orchestrates the entire testing process:

\begin{lstlisting}[caption=Main Test Function Structure]
template <typename hashtype>
bool LSHCollisionTest(const HashInfo * hinfo, bool extra, flags_t flags) {
    // Initialize output file
    std::ofstream out_file("collisionResults_" + std::string(hinfo->name) + "_" + ".csv");
    
    // Test multiple key sizes
    result &= LSHCollisionTestImpl<hashtype>(hinfo, 32, seed, flags, out_file);
    result &= LSHCollisionTestImpl<hashtype>(hinfo, 48, seed, flags, out_file);
    result &= LSHCollisionTestImpl<hashtype>(hinfo, 64, seed, flags, out_file);
    result &= LSHCollisionTestImpl<hashtype>(hinfo, 96, seed, flags, out_file);
    result &= LSHCollisionTestImpl<hashtype>(hinfo, 128, seed, flags, out_file);
    
    if (extra && !hinfo->isVerySlow()) {
        result &= LSHCollisionTestImpl<hashtype>(hinfo, 160, seed, flags, out_file);
        result &= LSHCollisionTestImpl<hashtype>(hinfo, 256, seed, flags, out_file);
    }
    
    return result;
}
\end{lstlisting}

\subsection{Implementation Algorithm}

The core algorithm in \texttt{LSHCollisionTestImpl} follows these steps:

\begin{enumerate}
    \item \textbf{Initialization}
    \begin{itemize}
        \item Calculate sequence length from key bits
        \item Set up similarity/dissimilarity vectors
        \item Initialize collision rate storage matrix
    \end{itemize}
    
    \item \textbf{For each dissimilarity level}
    \begin{itemize}
        \item Generate random sequences using \texttt{DataGeneration}
        \item Apply mutations using \texttt{MutationEngine}
        \item Calculate collision rates across hash family
    \end{itemize}
    
    \item \textbf{Data Output}
    \begin{itemize}
        \item Write collision rates to CSV file
        \item Format: rows = similarity levels, columns = trials
    \end{itemize}
\end{enumerate}

\begin{lstlisting}[caption=Core Algorithm Structure]
// For each dissimilarity index
for(dissimilarityIdx = 0; dissimilarityIdx < dissimrates.size(); dissimilarityIdx++) {
    
    // Generate sequences
    DataGeneration<hashtype> dataGen(&records, sequenceLength, keycount, dataGenSeed);
    dataGen.FillRecords(&records, dataGenSeed);
    
    // Apply mutations
    MutationEngine<hashtype> mutationEngine(sequenceLength, dissimrates[dissimilarityIdx], 
                                          MutationSeed, false, keycount);
    mutationEngine.MutateRecordsWithSubstitutions(&records, distances[dissimilarityIdx]);
    
    // Compute collision rates
    for(trialIdx = 0; trialIdx < keycount; trialIdx++) {
        float collisionCount = 0;
        
        // Test across hash family
        for (unsigned i = 0; i < hashcount; i++) {
            hashtype hashOrg, hashMut;
            seed_t currentSeed = currentHashSeed + (HashFamilySeedIncrement * i);
            
            hash(k1, keybytes, currentSeed, &hashOrg);
            hash(k2, keybytes, currentSeed, &hashMut);
            
            if (hashOrg == hashMut) {
                collisionCount++;
            }
        }
        
        collisionRate = collisionCount / static_cast<float>(hashcount);
        collisionRateVec[dissimilarityIdx][trialIdx] = collisionRate;
    }
}
\end{lstlisting}

\section{Data Generation and Mutation}

\subsection{Sequence Generation}

The \texttt{DataGeneration} class generates random DNA sequences:

\begin{itemize}
    \item Uses SMHasher3's \texttt{Rand} class for reproducible randomness
    \item Generates sequences using \texttt{SEQ\_DIST\_1} distribution
    \item Stores sequences in 2-bit encoded format for efficiency
\end{itemize}

\subsection{Mutation Engine}

The \texttt{MutationEngine} class introduces controlled mutations:

\begin{itemize}
    \item Implements substitution-based mutations only
    \item Uses Fisher-Yates shuffle for random position selection
    \item Ensures mutations result in different bases
    \item Maintains exact edit distance control
\end{itemize}

\begin{lstlisting}[caption=Mutation Algorithm]
// Fisher-Yates shuffle for position selection
for (int i = SequenceLength - 1; i > 0; i--) {
    int j = r.rand_range(i + 1);
    std::swap(indices[i], indices[j]);
}

// Apply mutations at selected positions
for(int i = 0; i < editDistance; i++) {
    int pos = indices[i];
    uint8_t original_base = extractBase(records, record_idx, pos);
    uint8_t new_base = mutate_base(original_base, rmutatebase);
    setBase(records, record_idx, pos, new_base);
}
\end{lstlisting}

\section{Hash Family Testing}

\subsection{Seed Generation Strategy}

The test employs a sophisticated seeding strategy to ensure independent hash functions:

\begin{align}
\text{Hash Family Seed} &= \text{Base Seed} + (0\text{x}9\text{e}3779\text{b}9 \times \text{Trial Index}) \\
\text{Individual Hash Seed} &= \text{Family Seed} + (0\text{x}9\text{e}3779\text{b}9 \times \text{Hash Index})
\end{align}

This ensures that:
\begin{itemize}
    \item Each trial uses a different family of hash functions
    \item Each hash within a family is independent
    \item Results are reproducible with fixed seeds
\end{itemize}

\subsection{Collision Detection}

For each sequence pair $(S_{original}, S_{mutated})$ and hash family $\{h_1, h_2, \ldots, h_k\}$:

\begin{equation}
\text{Collision Rate} = \frac{1}{k} \sum_{i=1}^{k} [h_i(S_{original}) = h_i(S_{mutated})]
\end{equation}

\section{Output Format and Analysis}

\subsection{CSV Output Structure}

The test generates a CSV file with the following structure:

\begin{center}
\begin{tabular}{|c|c|c|c|c|}
\hline
\textbf{Similarity Level} & \textbf{Trial 1} & \textbf{Trial 2} & \textbf{...} & \textbf{Trial 1024} \\
\hline
0.0 (100\% dissimilar) & rate$_{0,1}$ & rate$_{0,2}$ & ... & rate$_{0,1024}$ \\
0.1 (90\% dissimilar) & rate$_{1,1}$ & rate$_{1,2}$ & ... & rate$_{1,1024}$ \\
... & ... & ... & ... & ... \\
1.0 (0\% dissimilar) & rate$_{10,1}$ & rate$_{10,2}$ & ... & rate$_{10,1024}$ \\
\hline
\end{tabular}
\end{center}

\subsection{Expected Results}

For an ideal locality-sensitive hash function:

\begin{itemize}
    \item \textbf{High Similarity} (similarity $\rightarrow$ 1.0): Collision rate $\rightarrow$ 1.0
    \item \textbf{Low Similarity} (similarity $\rightarrow$ 0.0): Collision rate $\rightarrow$ $\frac{1}{2^{\text{hash bits}}}$
    \item \textbf{Monotonic Relationship}: Collision rate should increase with similarity
\end{itemize}

\section{Performance Considerations}

\subsection{Computational Complexity}

The total computational complexity is:
\begin{equation}
O(S \times K \times H \times L)
\end{equation}

Where:
\begin{itemize}
    \item $S$ = Number of similarity levels (11)
    \item $K$ = Number of key pairs (1024)
    \item $H$ = Number of hash functions (1024)  
    \item $L$ = Sequence length (varies by key size)
\end{itemize}

\subsection{Memory Usage}

Memory requirements include:
\begin{itemize}
    \item Sequence storage: $2 \times K \times \lceil L/4 \rceil$ bytes
    \item Collision rate matrix: $S \times K \times 4$ bytes
    \item Working memory for hash computations
\end{itemize}

\section{Error Handling and Validation}

\subsection{Bounds Checking}

The implementation includes comprehensive bounds checking:

\begin{lstlisting}[caption=Bounds Validation]
// Validate array access bounds
if (trialIdx * keybytes >= records.SeqTwoBitOrg.size() || 
    trialIdx * keybytes >= records.SeqTwoBitMut.size()) {
    printf("ERROR: Index out of bounds! trialIdx=%u, keybytes=%u\n", 
           trialIdx, keybytes);
    return false;
}
\end{lstlisting}

\subsection{Input Validation}

Key parameters are validated during initialization:
\begin{itemize}
    \item Sequence length $> 0$ and $> 13$ (minimum for meaningful analysis)
    \item Key count $> 0$
    \item Error rates in range $[0.0, 1.0]$
    \item Edit distances $\leq$ sequence length
\end{itemize}

\section{Integration with SMHasher3}

\subsection{Template Instantiation}

The test is instantiated for all supported hash types using the SMHasher3 instantiation system:

\begin{lstlisting}[caption=Template Instantiation]
INSTANTIATE(LSHCollisionTest, HASHTYPELIST);
\end{lstlisting}

This ensures compatibility with:
\begin{itemize}
    \item 32-bit hash functions
    \item 64-bit hash functions  
    \item 128-bit hash functions
    \item Variable-length hash functions
\end{itemize}

\subsection{Command Line Integration}

The test is accessible via:
\begin{lstlisting}[language=bash]
./SMHasher3 --test=LSHCollision [hash_name]
./SMHasher3 --test=LSHCollision --extra [hash_name]  # Extended testing
\end{lstlisting}

\section{Conclusion}

The LSH Collision Test provides a comprehensive framework for evaluating the locality-sensitive properties of hash functions on biological sequence data. By systematically varying sequence similarity and measuring collision rates across large hash families, it enables quantitative assessment of a hash function's suitability for applications requiring locality preservation.

The test's rigorous statistical approach, comprehensive parameter coverage, and integration with the SMHasher3 framework make it a valuable tool for hash function analysis in bioinformatics and related domains.

\end{document}